\section{Estimate Expected Value Without Fabricating Probabilities}

Expected value analysis is mathematically elegant when outcome distributions and probabilities can be calibrated with empirical data. However, at the earliest inception stage, a founder lacks empirical base rates. \textbf{Assigning arbitrary 70\% or 85\% probability weights to speculative spreadsheets creates a dangerous illusion of mathematical rigor.}

Instead of fabricating probabilities, apply 3 robust reasoning tools:
\begin{enumerate}[leftmargin=12pt,itemsep=2pt,topsep=1pt,parsep=0pt]
  \item \textbf{Bounded Scenario Analysis:} Explicitly model the Worst-Case Downside (zero sales after ₹50,000 spent), the Base Realizable Case (4 paying customers covering costs), and the High-Traction Case (12 customers validating the model).
  \item \textbf{Evidence-to-Action Decision Rules:} Determine in advance what exact empirical threshold must be achieved to trigger further investment (e.g., ``If at least 5 out of 20 interviewed engineers place a refundable deposit, we will build the 8-week curriculum; otherwise, we abandon the cohort hypothesis'').
  \item \textbf{Real Option Valuation:} Design the initial test as an affordable, bounded call option that purchases critical strategic information without committing the founder's entire net worth.
\end{enumerate}

\section{The 2-Stage Opportunity Selection Architecture}

Never compare fully realized, 5-year mature business models as if they currently exist. \textbf{Structure opportunity selection as a 2-stage discovery process}:

\begin{itemize}[leftmargin=12pt,itemsep=2pt,topsep=1pt]
  \item \textbf{Stage 1 (Experiment Selection):} Decide which 2 candidate opportunities are worth testing with low-cost, 30-day behavioral experiments.
  \item \textbf{Stage 2 (Business Commitment):} Select the core commercial vehicle only after empirical customer conversion and delivery data has been collected.
\end{itemize}

At the starting line, the correct executive decision might be: \textit{Execute 1 D2C diagnostic pilot in November, and 1 college curriculum pilot in January; then compare verified customer pull and unit economics.} It does not require prematurely crowning 1 business model as the permanent identity of the enterprise.

\section{Information Velocity versus Sales Velocity}

The fastest path to customer data is not always a rapid digital transaction:
\begin{itemize}[leftmargin=12pt,itemsep=1pt,topsep=1pt]
  \item A consumer landing page can generate 100 email signups in 48 hours, yet tell you zero about whether working engineers will complete a demanding 8-week curriculum.
  \item An enterprise partnership discussion may require 6 weeks to arrange, but a 45-minute discovery interview with a chief technology officer can immediately reveal that corporate budgets for AI testing are frozen until Q3.
\end{itemize}

\textbf{Select discovery experiments that match the natural decision clock of the customer.} Do not demand a 7-day sales cycle from an accredited university, and do not treat a free webinar signup as proof of customer commitment.

\section{The 6-Phase Research Sequence for Founders}

A rigorous market investigation unfolds across 6 structured phases:

\begin{enumerate}[leftmargin=14pt,itemsep=3pt,topsep=2pt]
  \item \textbf{Phase A (Framing \& Filtering):} Define 5 to 10 distinct buyer-job-trigger hypotheses. Pass them through the 6 non-negotiable gates. Discard non-compliant or capital-intensive options.
  \item \textbf{Phase B (Desk Research \& Base Rates):} Collect public government statistics, regulatory filings, competitor pricing models, and public job postings. Document sources and explicit dates.
  \item \textbf{Phase C (Qualitative Field Discovery):} Conduct 20 to 30 in-depth discovery interviews with recent buyers, non-buyers, and budget holders. Reconstruct their exact purchasing chronologies and workarounds.
  \item \textbf{Phase D (Unit Economic Modeling):} Model true direct delivery costs, customer acquisition effort, cash collection lags, and founder hours required per transaction.
  \item \textbf{Phase E (Behavioral Commitment Test):} Present a concrete, bounded commercial offer (a paid diagnostic clinic, a priced pilot proposal) to test authentic willingness to pay.
  \item \textbf{Phase F (Executive Decision):} Formally commit, narrow, pivot, or kill the opportunity based upon verified behavioral evidence.
\end{enumerate}

\section{The Research Triangulation Ladder}

Never rely upon a single stream of information. Reconcile 3 independent data channels:
\begin{itemize}[leftmargin=12pt,itemsep=2pt,topsep=1pt]
  \item \textbf{Macro Administrative Data:} Official higher-education statistics, government labor census reports, and corporate financial disclosures. Excellent for establishing macro population bounds; weak for predicting niche behavior.
  \item \textbf{Behavioral Field Evidence:} Qualitative discovery interviews, observed manual routines, and verified prepayments. Essential for understanding customer motivation; prone to small-sample selection bias.
  \item \textbf{Economic Mechanism Logic:} An explicit causal explanation of why the customer switches, why incumbents cannot retaliate, and how the startup captures surplus.
\end{itemize}

\textbf{When these 3 data streams disagree, you have located the boundary of your knowledge.} If macro reports predict booming corporate AI spending while 15 engineering directors report zero training budget, believe the directors.
